%\input{tcilatex}


\documentclass[11pt]{article}
\usepackage{amsfonts}

%%%%%%%%%%%%%%%%%%%%%%%%%%%%%%%%%%%%%%%%%%%%%%%%%%%%%%%%%%%%%%%%%%%%%%%%%%%%%%%%%%%%%%%%%%%%%%%%%%%%%%%%%%%%%%%%%%%%%%%%%%%%%%%%%%%%%%%%%%%%%%%%%%%%%%%%%%%%%%%%%%%%%%%%%%%%%%%%%%%%%%%%%%%%%%%%%%%%%%%%%%%%%%%%%%%%%%%%%%%%%%%%%%%
\usepackage[T1]{fontenc}
\usepackage[utf8]{inputenc}
\usepackage{lmodern}
\usepackage{amsmath,amssymb,amsthm,mathtools}
\usepackage[margin=1in]{geometry}
\usepackage{booktabs}
\usepackage{enumitem}
\usepackage{algorithm}
\usepackage{algpseudocode}
\usepackage{hyperref}
\usepackage{xcolor}

\setcounter{MaxMatrixCols}{10}
%TCIDATA{OutputFilter=Latex.dll}
%TCIDATA{Version=5.50.0.2953}
%TCIDATA{<META NAME="SaveForMode" CONTENT="1">}
%TCIDATA{BibliographyScheme=Manual}
%TCIDATA{LastRevised=Saturday, March 28, 2026 16:15:39}
%TCIDATA{<META NAME="GraphicsSave" CONTENT="32">}

%\input{tcilatex}

\begin{document}

\title{Life-cycle model with search frictions and skills}
\author{}
\date{\today}
\maketitle

\section{Model}

An individual is characterized by assets $a\geq 0$, employment status
$\ell \in \{0,1\}$, skill level
$g\in \{\bar{g}_{1},\ldots ,\bar{g}_{n_{g}}\}$, age
$j\in \{1,\ldots ,J\}$, and a fixed education type
$i\in \mathcal{I}\equiv \{\text{low},\text{high}\}$. Here $\ell =0$ denotes
unemployment and $\ell =1$ denotes employment. Retirement begins at age
$J_{R}$.

Education affects the environment through two type-specific objects:
$\phi_i>0$, an income shifter, and $\sigma_i\in (0,1)$, the job-separation
probability while employed. We assume $\sigma_{\text{low}}>
\sigma_{\text{high}}$, so higher-education individuals face a lower risk of
job loss.

At each age, the individual chooses next-period assets $a^{\prime }\geq 0$
and search effort $s\geq 0$. Search effort matters only while the agent is
unemployed and of working age, that is, when $\ell =0$ and $j<J_{R}$.

For all ages $j\leq J$, the Bellman equation is
\begin{equation*}
V_{i}(a,\ell ,g,j)=\max_{a^{\prime }\geq 0,\,s\geq 0}\left\{
u(c)-\kappa (s)+\beta \sum_{\ell ^{\prime }\in \{0,1\}}
\sum_{g^{\prime }\in \mathcal{G}}V_{i}(a^{\prime },\ell ^{\prime },g^{\prime
},j+1)\,\pi _{\ell }(\ell ,\ell ^{\prime }\mid s,i)\,\pi
_{g}(g,g^{\prime }\mid \ell ,j)\right\} ,
\end{equation*}%
where $\mathcal{G}=\{\bar{g}_{1},\ldots ,\bar{g}_{n_g}\}$, with terminal condition
\begin{equation*}
V_{i}(a,\ell ,g,J+1)=0.
\end{equation*}
Associated policy functions for next-period assets and search effort are
\begin{equation*}
Pol_{a^{\prime }}(a,\ell ,g,j,i), \qquad Pol_{s}(a,\ell ,g,j,i).
\end{equation*}

The period budget constraint is
\begin{equation*}
\left\{
\begin{array}{ll}
c+a^{\prime }=(1+r)a+\phi_i\left[w\,g\,\mathbb{I}_{\{\ell =1\}}+b\,\mathbb{I}
_{\{\ell =0\}}\right], & j<J_{R}, \\ 
c+a^{\prime }=(1+r)a+\phi_i\,pen, & j\geq J_{R}.%
\end{array}%
\right.
\end{equation*}%
Thus, at working ages labor income equals $\phi_iwg$ when employed and
$\phi_i b$ when unemployed, while retirees receive pension income
$\phi_i pen$. Feasible choices must satisfy $c>0$.

Period utility is
\begin{equation*}
u(c)=\frac{c^{1-\gamma}}{1-\gamma}, \qquad \kappa(s)=B_s s.
\end{equation*}
Since $\kappa(s)$ is strictly increasing, the optimal choice in all states in
which search brings no benefit is the lowest feasible effort, namely $s=0$.
In particular, search is irrelevant while employed and after retirement, when
employment status no longer affects payoffs. Likewise, after retirement the
payoff function no longer depends on $\ell$ or $g$, but we still keep those
variables in the state vector to maintain a uniform recursive formulation and
to match the Matlab implementation.

Employment transitions at working ages depend on search effort:
\begin{equation*}
\pi _{\ell }(\ell ,\ell ^{\prime }\mid s,i)=%
\begin{bmatrix}
& \ell ^{\prime }=0 & \ell ^{\prime }=1 \\ 
\ell =0 & 1-\pi ^{s}(s) & \pi ^{s}(s) \\ 
\ell =1 & \sigma _i & 1-\sigma _i%
\end{bmatrix}%
.
\end{equation*}%
An unemployed worker finds a job with probability $\pi ^{s}(s)$, while an
employed worker loses the job with probability $\sigma_i$. This matrix is
well defined provided $\pi^s(s)\in [0,1]$.

Skill transitions depend on current employment status. A simple three-state
illustration is
\begin{equation*}
\pi _{g}(g,g^{\prime }\mid \ell =0,j)=%
\begin{bmatrix}
& g^{\prime }=\bar{g}_{1} & g^{\prime }=\bar{g}_{2} & g^{\prime }=\bar{g}_{3}
\\ 
g=\bar{g}_{1} & 1 & 0 & 0 \\ 
g=\bar{g}_{2} & 1-p_{0} & p_{0} & 0 \\ 
g=\bar{g}_{3} & 0 & 1-p_{0} & p_{0}%
\end{bmatrix}%
,
\end{equation*}%
and
\begin{equation*}
\pi _{g}(g,g^{\prime }\mid \ell =1,j)=%
\begin{bmatrix}
& g^{\prime }=\bar{g}_{1} & g^{\prime }=\bar{g}_{2} & g^{\prime }=\bar{g}_{3}
\\ 
g=\bar{g}_{1} & p_{1}^{j} & 1-p_{1}^{j} & 0 \\ 
g=\bar{g}_{2} & 0 & p_{1}^{j} & 1-p_{1}^{j} \\ 
g=\bar{g}_{3} & 0 & 0 & 1%
\end{bmatrix}%
.
\end{equation*}%
Under this ladder structure, unemployment weakly depreciates skills, while
employment weakly improves them. To make the matrices stochastic, one needs
$p_0\in [0,1]$ and $p_1^j\in [0,1]$ for every age $j$. The age index in
$p_{1}^{j}$ allows skill accumulation to vary over the life cycle, for
example by making it stronger at younger ages. 

\subsection{Distribution}

Let $\varpi_i$ denote the fixed population share of permanent education type
$i\in \mathcal{I}$, with $\sum_{i\in \mathcal{I}}\varpi_i=1$. Let
$m_j\geq 0$ denote the cross-sectional age-mass profile, with
$\sum_{j=1}^{J}m_j=1$. In the code this object is stored as
$m_j=Params.mewj(j)$. The current application sets $m_j=1/J$ for all $j$,
but the formulation allows for a general age distribution, for example one
generated by survival risk. Let $\tilde{\mu}(a,\ell,g,j,i)$ denote the
distribution over $(a,\ell,g,i)$ within age $j$, before applying the
cross-sectional age weights, and let $\mu(a,\ell,g,j,i)$ denote the joint
cross-sectional distribution. These objects are related by
\begin{equation*}
\mu(a,\ell,g,j,i)=m_j\,\tilde{\mu}(a,\ell,g,j,i).
\end{equation*}
Because permanent type is fixed over the life cycle, both distributions keep
track of $i$ explicitly rather than integrating it out. After retirement, the
distribution still carries $\ell$ and $g$, even though they are
payoff-irrelevant, in order to preserve the same state-space representation
at all ages.

Given these policy functions, the within-age distribution evolves according
to
\begin{eqnarray*}
\tilde{\mu}(a^{\prime},\ell^{\prime},g^{\prime},j+1,i)
&=& \sum_{a,\ell,g}\tilde{\mu}(a,\ell,g,j,i)\,
\mathbb{I}\left\{a^{\prime}=Pol_{a^{\prime}}(a,\ell,g,j,i)\right\} \\
&&\times \pi_{\ell}(\ell,\ell^{\prime}\mid Pol_s(a,\ell,g,j,i),i)\,
\pi_g(g,g^{\prime}\mid \ell,j),
\end{eqnarray*}%
for $j=1,\ldots,J-1$.
Equivalently, the cross-sectional distribution satisfies
\begin{eqnarray*}
\mu(a^{\prime},\ell^{\prime},g^{\prime},j+1,i)
&=& \frac{m_{j+1}}{m_j}\sum_{a,\ell,g}\mu(a,\ell,g,j,i)\,
\mathbb{I}\left\{a^{\prime}=Pol_{a^{\prime}}(a,\ell,g,j,i)\right\} \\
&&\times \pi_{\ell}(\ell,\ell^{\prime}\mid Pol_s(a,\ell,g,j,i),i)\,
\pi_g(g,g^{\prime}\mid \ell,j),
\end{eqnarray*}%
whenever $m_j>0$.

Let $g_0\in \mathcal{G}$ denote the median skill level. The initial
within-age distribution at age $j=1$ is given by
\begin{equation*}
\tilde{\mu}(a,\ell,g,1,i)=\varpi_i\,\mathbb{I}\left\{a=0\right\}
\mathbb{I}\left\{\ell=0\right\}\mathbb{I}\left\{g=g_0\right\},
\end{equation*}%
so all newborn agents start with zero assets, unemployed, and at the median
skill level $g_0$, while permanent types are distributed according to
$\varpi_i$. Therefore the age-$1$ cross-sectional distribution is
\begin{equation*}
\mu(a,\ell,g,1,i)=m_1\,\varpi_i\,\mathbb{I}\left\{a=0\right\}
\mathbb{I}\left\{\ell=0\right\}\mathbb{I}\left\{g=g_0\right\}.
\end{equation*}

\subsection{Life-cycle profiles}

Consider a generic variable
$x=x(a,\ell,g,j,i)$ defined on the individual state space. Examples include
assets, search effort, and employment status. Using the cross-sectional
distribution $\mu(a,\ell,g,j,i)$, the average of $x$ conditional on age $j$
is
\begin{equation*}
\mathbb{E}[x\mid j]
=\frac{\sum_{i\in\mathcal{I}}\sum_{a,\ell,g}x(a,\ell,g,j,i)\,
\mu(a,\ell,g,j,i)}
{\sum_{i\in\mathcal{I}}\sum_{a,\ell,g}\mu(a,\ell,g,j,i)}.
\end{equation*}
Since $\mu$ is the joint cross-sectional distribution and the age mass at
age $j$ is $m_j$, the denominator equals $m_j$. Hence one can also write
\begin{equation*}
\mathbb{E}[x\mid j]
=\frac{1}{m_j}\sum_{i\in\mathcal{I}}\sum_{a,\ell,g}
x(a,\ell,g,j,i)\,\mu(a,\ell,g,j,i).
\end{equation*}

If instead one wants the average conditional on both age $j$ and permanent
type $i$, then
\begin{equation*}
\mathbb{E}[x\mid j,i]
=\frac{\sum_{a,\ell,g}x(a,\ell,g,j,i)\,\mu(a,\ell,g,j,i)}
{\sum_{a,\ell,g}\mu(a,\ell,g,j,i)}.
\end{equation*}
Because permanent type shares are fixed over the life cycle, the denominator
is $m_j\varpi_i$, so equivalently
\begin{equation*}
\mathbb{E}[x\mid j,i]
=\frac{1}{m_j\varpi_i}\sum_{a,\ell,g}
x(a,\ell,g,j,i)\,\mu(a,\ell,g,j,i).
\end{equation*}

If one wants the average of $x$ conditional on age $j$ and employment
status $\ell=0$, then
\begin{equation*}
\mathbb{E}[x\mid j,\ell=0]
:=\frac{\sum_{i\in\mathcal{I}}\sum_{a,g}x(a,0,g,j,i)\,\mu(a,0,g,j,i)}
{\sum_{i\in\mathcal{I}}\sum_{a,g}\mu(a,0,g,j,i)}.
\end{equation*}
Similarly, conditional on age $j$ and employment status $\ell=1$,
\begin{equation*}
\mathbb{E}[x\mid j,\ell=1]
:=\frac{\sum_{i\in\mathcal{I}}\sum_{a,g}x(a,1,g,j,i)\,\mu(a,1,g,j,i)}
{\sum_{i\in\mathcal{I}}\sum_{a,g}\mu(a,1,g,j,i)}.
\end{equation*}
More compactly, for any $\bar{\ell}\in\{0,1\}$,
\begin{equation*}
\mathbb{E}[x\mid j,\ell=\bar{\ell}]
:=\frac{\sum_{i\in\mathcal{I}}\sum_{a,g}x(a,\bar{\ell},g,j,i)\,\mu(a,\bar{\ell},g,j,i)}
{\sum_{i\in\mathcal{I}}\sum_{a,g}\mu(a,\bar{\ell},g,j,i)}.
\end{equation*}

Other moments by age, or jointly by age and permanent type, can be computed
in exactly the same way by replacing $x$ with the relevant function of the
state variables.

\paragraph{Implementation note.}

In the VFI Toolkit, this problem is implemented as a finite-horizon model
with semi-exogenous states and no purely exogenous Markov state $z$. We
define the semi-exogenous state as $semiz=(\ell ,g)$. The reason is that the
employment transition depends on the decision variable $d=s$, while skill
transitions depend on current employment status $\ell$. Hence $\ell$ is
directly semi-exogenous, and $g$ is naturally treated as a second
semi-exogenous state.

\begin{center}
\begin{tabular}{lll}
\toprule
Variable & Description & Toolkit name \\
\midrule
$s$ & Search effort & $d$ \\
$a^{\prime }$ & Next-period endogenous state & $aprime$ \\
$a$ & Current endogenous state & $a$ \\
$\ell$ & Employment state & $semiz1$ \\
$g$ & Skill level & $semiz2$ \\
\bottomrule
\end{tabular}
\end{center}

In that notation, the relevant objects are the semi-exogenous grid
$semiz\_grid\_J$ and the transition array $\pi _{semiz,J}$. In full
generality these can be entered as multidimensional arrays:
\begin{equation*}
semiz\_grid\_J(semiz,j,ptype), \qquad \pi _{semiz,J}(semiz,semiz^{\prime },d,j,ptype).
\end{equation*}
If the semi-exogenous grid is the same for all ages and permanent types,
however, it can be entered in stacked form as a column vector
$semiz\_grid\_J(semiz)$ rather than as a joint grid. See the VFI Toolkit
documentation or code for further details on this notation and data layout.

\end{document}
